%%%%%%%%%%%%%%%%%%%%%%%%%%%%%%%%%%%%%%%%%%%%%%%%%%%%%%%%%%%%%%%%%%%%%%
%                                                                    %
% Multiperfect numbers on lines of the Pascal triangle,              %
% by                                                                 %
% Florian Luca and Juan Luis Varona                                  %
%                                                                    %
% Journal of Approximation Theory                                    %
%                                                                    %
% DATE OF LAST MODIFICATION OF THIS FILE: November 26, 2008          %
%                                                                    %
% Written with Elsevier's elsart class                               %
%                                                                    %
%%%%%%%%%%%%%%%%%%%%%%%%%%%%%%%%%%%%%%%%%%%%%%%%%%%%%%%%%%%%%%%%%%%%%%

\documentclass[twoside,amsthm]{elsart}

%-------------------------------
\journal{J.\ Number Theory}
\volume{129}
%\issue{5}
\pubyear{2009}
%\copyear{2008}
\firstpage{1136}
\lastpage{1148}
%\company{Elsevier Science}
%\copyrightholder{Elsevier Science}

\runauthor{F. Luca and J. L. Varona}
\runtitle{Multiperfect numbers on lines of the Pascal triangle}
\AtBeginDocument{\pagestyle{myheadings}}
%-------------------------------

% \usepackage{amsmath,amssymb,amsthm}
\usepackage{amsmath,amssymb} % <- elsart loads amsthm

\usepackage{url}

\journal{Journal of Number Theory}

%\newtheorem{lem}{Lemma}
%\newtheorem{lemma}[lem]{Lemma}

\newtheorem{lemma}{Lemma} % <-- To have its own numbers

%\newtheorem{prop}{Proposition}
%\newtheorem{proposition}[prop]{Proposition}

%%% La clase elsart ya incluye \newtheorem{thm}{Theorem} en sus definiciones:
% \newtheorem{thm}{Theorem}
%\newtheorem{theorem}[thm]{Theorem}

%\newtheorem{cor}{Corollary}
%\newtheorem{corollary}[cor]{Corollary}

%\newtheorem{conj}{Conjecture}
%\newtheorem{conjecture}[conj]{Conjecture}

%\newtheorem{defi}{Definition}
%\newtheorem{definition}[defi]{Definition}

%\newtheorem{exam}{Example}
%\newtheorem{example}[exam]{Example}

%\newtheorem{prob}{Problem}
%\newtheorem{problem}[prob]{Problem}

%\newtheorem{ques}{Question}
%\newtheorem{question}[ques]{Question}

%% DEFINITIONS

\def\cA{\mathcal{A}}
\def\cB{\mathcal{B}}
%\def\cC{\mathcal{C}}
%\def\cE{\mathcal{E}}
%\def\cF{\mathcal{F}}
\def\cI{\mathcal{I}}
%\def\cL{\mathcal{L}}
%\def\cM{\mathcal{M}}
%\def\cN{\mathcal{N}}
%\def\cR{\mathcal{R}}
%\def\cS{\mathcal{S}}
\def\cP{\mathcal{P}}
\def\cQ{\mathcal{Q}}
%\def\cT{\mathcal{T}}
%\def\cX{\mathcal{X}}
\def\N{\mathbb{N}}
\def\Z{\mathbb{Z}}
%\def\R{\mathbb{R}}
%\def\C{\mathbb{C}}
%\def\eps{\varepsilon}


%--------------------------------------------------------------
%%%% CAMBIOS PARA ADAPTAR (aproximadamente) LA CLASE alsart QUE DISTRIBUYE Elsevier
%%%% CON LA QUE LUEGO REALMENTE UTILIZAN PARA PRODUCIR LAS VERSIONES FINALES DE
%%%% LOS ARTICULOS.
\makeatletter
%% Truco para poder mostrar las p\'aginas del art\'{\i}culo tal como
%% lo public\'o Elsevier:
% \def\@jou@vol@pag{\@journal\ \@volume\ (\the\@pubyear)\ \@pagerange} % <- ORIGINAL
\newcounter{publishedfirstpage}\setcounter{publishedfirstpage}{\thefirstpage}
\firstpage{1}
\newcounter{publishedlastpage}\setcounter{publishedlastpage}{\thelastpage}
\def\@jou@vol@pag{\@journal\ \@volume\ (\the\@pubyear)
  \thepublishedfirstpage--\thepublishedlastpage.}
%% MODIFICACIONES A LAS keywords y AL MSC
\def\@keywordheading{{\it Keywords:\ }}
\let\originalkeyword\keyword
\def\keyword{\originalkeyword
  \def\sep{\unskip; }
  \def\MSC{\par\leavevmode\hbox {\it MSC:\ }}%
  \ignorespaces}
%% ESTOS DOS PARAMETROS NO SE PARA QUE VALEN, PERO LOS HE CAMBIADO POR SI ACASO
%% (salvo el \@abstractsize que lo he dejado como estaba):
\def\@abstractsize{\fontsize{\@ixpt}{11pt}\selectfont} %% {\@ixpt}{11pt}
\def\@articletypesize{\fontsize{\@xpt}{10pt}\selectfont} %% {\@xiiipt}{13pt}
%% ESTOS TAMANOS {letra}{interlineado} LOS HE PUESTO MAS PEQUENOS
%% (a la derecha pongo lo que venia en la clase elsart por defecto):
\def\normalsize{\@setfontsize\normalsize\@xpt{12pt}% %% \@xiipt{14.5pt}
\abovedisplayskip 12\p@ \@plus 2\p@ \@minus 2\p@%
\belowdisplayskip \abovedisplayskip
\abovedisplayshortskip \z@ \@plus 2\p@%
\belowdisplayshortskip 3.5\p@ \@plus 2\p@ \@minus 2\p@
\let\@listi\@listI}
\def\small{\@setfontsize\small\@ixpt{10.5}% %% \@xipt{13.6}
\abovedisplayskip 11\p@ \@plus3\p@ minus6\p@
\belowdisplayskip \abovedisplayskip
\abovedisplayshortskip  \z@ \@plus3\p@
\belowdisplayshortskip  6.5\p@ \@plus3.5\p@ minus3\p@
\def\@listi{\leftmargin\leftmargini
 \parsep 2.5\p@ \@plus1\p@ minus0.5\p@ \itemsep \parsep %% 4.5   2   1
            \topsep 3\p@ \@plus1\p@ minus1\p@}}      %% 9   3   5
\let\footnotesize=\small
\let\@xviiipt\@xviipt
\def\scriptsize{\@setfontsize\scriptsize\@viipt{8.5}} %% \@viiipt{9.5}
\def\tiny{\@setfontsize\tiny\@vpt{6}} %% \@vipt{7}
\def\large{\@setfontsize\large\@xiipt{14}} %% \@xivpt{18}
\def\Large{\@setfontsize\Large\@xivpt{18}} %% \@xviipt{22}
\def\LARGE{\@setfontsize\LARGE\@xviipt{19}} %% \@xxpt{22}
\def\huge{\@setfontsize\huge\@xxpt{22}} %% \@xxvpt{27}
\let\Huge=\huge
%% ESTOS ESPACIOS {antes}{despues} LOS HE PUESTO MAS PEQUENOS:
\def\section{\@startsection{section}{1}{\z@}{0.8\@bls %% default elsart: 1.5\@bls
  \@plus .3\@bls \@minus .2\@bls}{0.4\@bls}{\normalsize\bfseries}}
\def\subsection{\@startsection{subsection}{2}{\z@}{0.6\@bls
  \@plus .2\@bls \@minus .1\@bls}{0.3\@bls}{\normalsize\itshape}}
\def\subsubsection{\@startsection{subsubsection}{3}{\z@}{0.4\@bls
  \@plus .1\@bls}{0.0001pt}{\normalsize\itshape}}
\def\paragraph{\@startsection{paragraph}{4}{\z@}{1.25ex \@plus %% default elsart: 3.25ex
  0.2ex \@minus 0.2ex}{-1em}{\normalsize\bfseries}}
\parskip 2\p@ \@plus 1\p@ %% DEFAULT elsart: \parskip 1pc \@plus 1\p@
\parindent 12pt           %% DEFAULT elsart: \parindent \z@
\def\Elproofname{Proof.}  %% DEFAULT elsart: PROOF.
  \def\th@plain{%
    \thm@preskip\medskipamount %% DEFAULT elsart: \parskip
    \thm@postskip\z@
    \itshape % body font
  }
%% ANALOGAMENTE PODRIA MODIFICAR \th@definition Y \th@remark
\makeatother

%%%% CAMBIOS PARA QUE PONGA "HEADINGS":
\makeatletter
% \ps@myheadings est\'a en article.cls:
\def\ps@myheadings{%
    \let\@oddfoot\@empty\let\@evenfoot\@empty
    \def\@evenhead{\thepage\hfil\small\slshape\leftmark\hfil}%
    \def\@oddhead{\hfil{\small\slshape\rightmark}\hfil\thepage}%
    \let\@mkboth\@gobbletwo
    \let\sectionmark\@gobble
    \let\subsectionmark\@gobble
    }
\def\leftmark{\@runauthor}
\def\rightmark{\@runtitle}
\headheight \baselineskip %% DEFAULT elsart: \z@
\headsep 20\p@            %% DEFAULT elsart: \z@
\makeatother
%--------------------------------------------------------------



%%%%%%%%%%%%%%%%%%%%%%%%%%%%%%%%%%%%%%%%%
\begin{document}
%%%%%%%%%%%%%%%%%%%%%%%%%%%%%%%%%%%%%%%%%

\begin{frontmatter}

\title{Multiperfect numbers on lines of the Pascal triangle}

\author{Florian Luca\thanksref{SEP}}
\ead{fluca@matmor.unam.mx}
\address{Instituto de Matem\'aticas,
Universidad Nacional Aut\'onoma de M\'exico,\\
C.P.~58089, Morelia, Michoac\'an, Mexico} % "Mexico" written in English
\thanks[SEP]{Supported in part by Grant SEP-CONACyT 79685.}

\author{Juan Luis Varona\thanksref{MTM2006}}
\ead{jvarona@unirioja.es}
% \ead[url]{http://www.unirioja.es/cu/jvarona/}
\address{Dpto. de Matem\'aticas y Computaci\'on,
Universidad de La Rioja,\\
Calle Luis de Ulloa s/n,
C.P.~26004, Logro\~no, Spain}
\thanks[MTM2006]{Supported by grant MTM2006-13000-C03-03 of the DGI.}

\date{Received 16 March 2008. Revised 27 October 2008}

%\maketitle

\begin{abstract}
A number $n$ is said to be multiperfect (or multiply perfect) if $n$
divides its sum of divisors $\sigma(n)$. In this paper, we study the
multiperfect numbers on straight lines through the Pascal triangle.
Except for the lines parallel to the edges, we show that all other
lines through the Pascal triangle contain at most finitely many
multiperfect numbers. We also study the distribution of the numbers
$\sigma(n)/n$ whenever the positive integer $n$ ranges through the
binomial coefficients on a fixed line through the Pascal triangle.
\end{abstract}

\begin{keyword}
Multiperfect \sep multiply perfect \sep perfect \sep
sum of divisors \sep
Pascal triangle \sep binomial coefficients \sep
Catalan numbers.
\smallskip
\MSC 11A25 \sep 11B65.
\end{keyword}

\end{frontmatter}


%%%%%%%%%%%%%%%%%%%%%%%%%%%%%%%%%%%%%%%%%
\section{Introduction and main results}
%%%%%%%%%%%%%%%%%%%%%%%%%%%%%%%%%%%%%%%%%

For a positive integer $n$ we put $\sigma(n)$ for the sum of its
divisors. Given an integer $k$, the number $n$ is said to be
multiperfect, multiply perfect, or $k$-fold perfect if $\sigma(n) =
kn$. Of course, ordinary perfect numbers are $2$-fold perfect. The
single $1$-fold perfect is the trivial case $n=1$. The $3$-fold
perfect numbers are also called triperfect, and only six of them are
known: they are $120$, $672$, $523776$, $459818240$, $1476304896$, 
$51001180160$. All of them were already known by the seventeenth
century. Several multiperfect numbers are also known for every $k\le
11$. Their number varies from thousands for $k=8,9,10$, to only one
for $k=11$ which has more than a thousand decimal digits and was
discovered in 2001. Descartes discovered the first $4$-fold number,
and Fermat the first $5$-fold number, respectively. Dickson's
\textit{History of the Theory of Numbers} \cite[p.~33--38]{D}
records a long interest in such numbers. See
also~\cite[Section~B2]{Guy}, or the web page~\cite{njas} for more
details and references.

Except for the well-known Euclid-Euler rule for $k=2$, no formula to
generate multiperfect numbers is known.  Lehmer~\cite{Leh} proved
that if $n$ is odd, then $n$ is perfect just if $2n$ is $3$-fold
perfect. Moreover, no odd multiperfect number is known. There are
several conjectures on the size of $k$ in what relates to the size
of~$n$. For example, from the maximal order of the sum of divisors
function, it is known that there exists a positive constant $c$ such
that the inequality $\sigma(n)/n>c\log\log n$ holds for infinitely
many positive integers $n$, where here and from now on we use $\log
$ for the natural logarithm. Contrary to this inequality, Erd\H{o}s
conjectured that if there were infinitely many multiperfect numbers,
then $k = o(\log\log n)$ as $n\to\infty$ through multiperfect
numbers. It has even been suggested there may be only finitely many
$k$-fold perfect numbers altogether with $k \ge 3$, and it is
further believed that all multiperfect numbers with $k=3$, $4$, $5$,
$6$, and $7$ are known.

There are several results in the literature addressing perfect and
multiperfect numbers of various shapes. For example, Pomerance
\cite{Po} proposed as a problem to find all positive integers $n$ such
that $n!$ is multiperfect. In the solution \cite{Ev} to the above
problem, it is shown that this happens only for $n=1$, $3$,~$5$. In
\cite{Lu1}, it is shown that there is no Fibonacci number which is
perfect and in \cite{Lu3} it was shown that there are at most
finitely many Fibonacci numbers which are multiperfect.
In~\cite{Lu2}, it is shown that no Fermat number; i.e., number of
the form $2^{2^n}+1$ for some non-negative integer $n$, is perfect,
and the method of proof shows easily that such numbers are not
multiperfect either.

There are various extensions of perfect and multiperfect numbers in
the literature. For example, a number $n$ is called superperfect if
$\sigma(\sigma(n))=2n$. More generally, a number $n$ is called
$(m,k)$-perfect if $\sigma^{(m)}(n)=kn$, where $\sigma^{(m)}$
denotes the $m$th fold iterate of the sum of divisors function.

It has been shown that there are only finitely many multiperfect
numbers of the form $\binom{2n}{n}$ (see~\cite{L-P, L-S}). In this
paper, we take a closer look at this problem and we study the
multiperfect as well as $(2,k)$-perfect numbers (for any positive
integer~$k$) on straight lines on the Pascal triangle.
Our theorem here is the following:

\begin{thm}
\label{thm:main} Let $L$ be a line passing through the Pascal
triangle and not parallel to any of its sides. Then the following
hold:
\begin{itemize}\leftskip=10pt\labelsep=6pt
\item[\normalfont (i)] There are only finitely many binomial 
coefficients $x$ belonging to $L$ such that $x\mid \sigma(x)$.
\item[\normalfont (ii)] There are only finitely many binomial 
coefficients $x$ belonging to $L$ such that 
$x\mid \sigma(\sigma(x))$.
\end{itemize}
\end{thm}

Our theorem above does not say anything about the lines parallel to
the edges; they are $\{\binom{n}{m}:n\ge m \}$ with a fixed value of
$m$ or their symmetries with respect to the middle axis
$\{\binom{n}{n-m}:n\ge m\}$. When $m=1$, we have $\binom{n}{1} = n$,
so saying something about these numbers amounts to saying something
about all the numbers which are multiperfect. When $m=2$, we have
$\binom{n}{2} = \frac{n(n-1)}{2}$. All even perfect numbers are of
this form with $n=2^r$, where $2^r-1$ is a prime number. However,
for $m \ge 3$, and limiting ourselves to perfect numbers only, we
have the following result.

\begin{thm}
\label{thm:kge3} There is no perfect number of the form
$\binom{n}{m}$ with $n\ge 2m\ge 6$.
\end{thm}

The proof of Theorem~\ref{thm:main} is effective and relies on the
Prime Number Theorem and the Primitive Divisor Theorem for the
Mersenne numbers, which have the form $M_n=2^n-1$ for all $n\ge 1$.
That is, given the line $L$, one can write down an explicit upper
bound for the size of the largest possible multiperfect number on~$L$.
Below we give a result identifying all the multiperfect numbers
as well as the $(2,k)$-perfect numbers (for any positive integer~$k$) 
on the central line of the Pascal triangle. For a positive
integer $n$ let $B_n=\binom{2n}{n}$ be the $n$th middle binomial
coefficient and $C_n=\frac{1}{n+1}B_n$ be the $n$th Catalan number.


\begin{thm}
\label{thm:2nn}
\begin{itemize}\item[]\end{itemize} % <--- Dirty trick to change the line.
\begin{itemize}\leftskip=10pt\labelsep=6pt
\item[\normalfont (i)]
The only $n$ such that $B_n$ is multiperfect is $n=2$. The only $n$
such that $C_n$ is multiperfect is $n=1$.
\item[\normalfont (ii)] The only
$n$ such that $B_n\mid \sigma(\sigma(B_n))$ is $n=1$. Furthermore,
the only $n$ such that $C_n\mid\sigma(\sigma(C_n))$ are $n=1,~2$ 
and~$5$.
\end{itemize}
\end{thm}

The occurrence of members of specific sequences on lines in the
Pascal triangle is a question that has been studied before. For
example, in \cite{LuPa} it is shown that if $(u_n)_{n\ge 0}$ is any
fixed binary recurrent sequence of integers and $L$ is any line of
the Pascal triangle, except for an edge or a line of the form
$\{\binom{n}{1}:n\ge 1\}$ or its symmetric $\{\binom{n}{n-1}:n\ge
1\}$, there are at most finitely many members of the sequence
$(u_n)_{n\ge 0}$ belonging to~$L$. The problem~\cite{Lu4P} asks to
show that the largest Fibonacci number which is also a Catalan
number is~$5$.

Divisor sums of binomial coefficients have been studied also in the
recent paper \cite{Lu6} in the following context. It is known that
the average value of $\sigma(n)/n$ when $n$ ranges in the interval
$[1,x]$ approaches $\pi^2/6$ as $x$ tends to infinity. When
$(u_n)_{n\ge 0}$ is a nondegenerate linearly recurrent sequence of
integers, then the average value of $\sigma(|u_n|)/|u_n|$ when $n$
ranges in the interval $[1,x]$  through values such that $u_n\ne 0$
approaches a finite limit when $x\to\infty$. This is a result of
Shparlinski from \cite{Shp} (see also~\cite{LuSh}). For binomial
coefficients, the situation is different. Indeed, it is shown
in~\cite{Lu6} that there exists a positive constant $c$ such that the
inequality
$$
\frac{1}{n+1}\sum_{0\le m\le n}
\frac{\sigma\left(\binom{n}{m}\right)}{\binom{n}{m}}
> c\log\log\log n
$$
holds for all sufficiently large $n$ and further that
$$
\frac{1}{2^n} \sum_{0\le m\le n} \sigma\left(\binom{n}{m}\right)
\to \infty
\quad {\text{as}}~n\to\infty.
$$
The above estimates should be consistent with our intuition that
most binomial coefficients are divisible by many small primes,
therefore their divisor sums should be quite a bit larger than the
numbers themselves.

In the present paper, we take again a line $L$ passing through the
Pascal triangle and study the distribution of the numbers
$\sigma(x)/x$, when $x$ ranges through members of~$L$. For the
purpose of the next result, we shall assume that $L$ contains
infinitely many members of the Pascal triangle. The next result
shows that numbers $\sigma(x)/x$ tend to behave differently when $x$
ranges through a line $L$ nonparallel to the sides than when $x$
ranges through one line that is parallel to the sides. Informally
speaking, $\sigma(x)/x$ tends to infinity when $L$ is not parallel
to the sides, while these numbers form a dense set in $[1,\infty)$
when the line is parallel to the sides.

\begin{thm}
\label{thm:dens} Let $L$ be a fixed line through the Pascal triangle
containing infinitely many binomial coefficients and which is not
one of the two edges.
\begin{itemize}\leftskip=10pt\labelsep=6pt
\item[\normalfont (i)] If $L$ is not parallel to any of the edges, 
then there exists a positive constant $c_L$ depending on the line such 
that the inequality $\sigma(x)/x>(\log\log x)^{c_L}$ holds for all 
binomial coefficients $x$ on $L$ which are sufficiently large.
\item[\normalfont (ii)] If $L=\{\binom{n}{m}~:~n\ge m\}$, where 
$m\ge 1$ is a fixed positive integer, then $\{\sigma(x)/x:x\in L\}$ 
is dense in $[1,\infty)$.
\end{itemize}
\end{thm}

Density questions for the set $\{\sigma(u_n)/u_n:n\ge 1\}$, where
$(u_n)_{n\ge 1}$ is some interesting sequence of positive integers
have been studied before. For example, it was shown in~\cite{Lu5}
that $\{\sigma(M_n)/M_n:n\ge 1\}$ is dense in $[1,\infty)$, where
again $M_n=2^n-1$ is the $n$th Mersenne number. The proof of this
result applies also when the sequence of Mersenne numbers is
replaced by the sequence of Fibonacci numbers. In \cite{LuPo}, it is
shown that $\{\sigma(\phi(n))/n : n\ge 1\}$ is dense in
$[\frac{1}{2}+\frac{1}{2^{34}-1},\infty)$, where $\phi(n)$ is the
Euler totient function. It is not very hard to
prove that $\{\sigma(\sigma(n))/n:n\ge 1\}$ is dense in
$[1,\infty)$, but it is not known whether
$\{\sigma(\sigma(\sigma(n)))/n:n\ge 1\}$ is dense in $[1,\infty)$.

As usual, we will use the Landau symbols $O$ and $o$ as well as the
Vinogradov symbols $\gg$, $\ll$ and $\asymp$ with their regular
meaning. Recall that the statements $A=O(B),~A\ll B$ and $B\gg A$
are all equivalent to the fact that $|A|<cB$ holds for some positive
constant~$c$, $A\asymp B$ means that both $A\ll B$ and $B\ll A$
hold, and $A=o(B)$ means that $A/B$ approaches zero. We use $\pi(t)$
for the number of primes $p\le t$.



%%%%%%%%%%%%%%%%%%%%%%%%%%%%%%%%%%%%%%%%%
\section{Proof of Theorem~\ref{thm:main}}
%%%%%%%%%%%%%%%%%%%%%%%%%%%%%%%%%%%%%%%%%

Let us start by presenting a parametrization of the lines in the
Pascal triangle.

\begin{lemma}
\label{lem:param}
Let $L \subseteq \{(n,m):n\ge m\}$ be a line of the Pascal triangle
not parallel to any of its sides and containing infinitely
many elements. Then, $L$ can be parametrized as
\begin{equation}
\label{eq:system}
\begin{cases}
n=at+b,\\
m=ct+d,
\end{cases}
\qquad t \in \Z,\quad t \ge t_0,
\end{equation}
for some integers $a,b,c,d$ with $a$ and $c$ coprime and $0< c < a$.
\end{lemma}

\begin{pf}
If the line $L$ is parallel to one of the sides,
then either $m$ is a constant or $n-m$ is a constant.
Otherwise, there exist rational numbers $A,~B$ such that $n=Am+B$.
Writing $A=a/c$ with coprime integers $a$ and $c>0$, and writing
$m=ct+r$, where $r\in \{0,1,\ldots,c-1\}$ and $t$ is an integer, we
get that
$$
n=at+\left(\frac{ar}{c}+\frac{e}{f}\right),
$$
where we put $B=e/f$ with coprime integers $e$ and $f>0$. Since $r$
can take only finitely many values, we may assume that $r$ is fixed.
Since $ar/c+e/f$ is an integer, we get that $f\mid c$, therefore
$c=fc_1$ holds with some positive integer $c_1$, and further
$r=c_1r_1$ holds with some nonzero integer~$r_1$. Furthermore,
$(ar_1+e)/f$ is an integer. Thus, putting $b:=(ar_1+e)/f$ and
$d:=r$, we get~\eqref{eq:system}.
Since $c>0$ and $m$ is positive, it follows that $t>0$ if $m$ is
sufficiently large. The first relation above together with the fact
that $n\ge m$ can be arbitrarily large implies first that $a \ge 1$.
Next, again since $n\ge m$, it follows by looking at large values of
$t$ that $0< c \le a$.
Finally, we can assume that $0<c<a$, otherwise if $a=c$ we
get that $n-m$ remains constant.
\end{pf}

Let us now prove Theorem \ref{thm:main}. Let $L \subseteq
\{(n,m) : n\ge m\}$ be the line of the Pascal triangle; by Lemma
\ref{lem:param} above, it can be parametrized as shown
in~\eqref{eq:system}. Furthermore, by the symmetry $\binom{n}{m} =
\binom{n}{n-m}$, and by replacing $m$ with $n-m$ if needed, it
follows that we may assume that $0<c\le a/2$.

Now let us look at the binomial coefficient:
$$
x := \binom{n}{m} = \binom{at+b}{ct+d} = \frac{(at+b)(at+b-1) \cdots
((a-c)t+b-d+1)}{(ct+d)!}.
$$
Since $a-c\ge c$, it follows that for large $t$, all primes $p$ in
the interval $\cI = [(a-c)t+b-d+1, at+b]$ divide $x$ with at
most $O(1)$ exceptions. These exceptional values can occur only when
$a-c=c$ (i.e., $a=2c$, therefore $c=1$ and $a=2$, since we have
already assumed that $a$ and $c$ are coprime), and $b-d+1<d$, in
which case we can take the constant implied by the above $O(1)$ to
be $\max\{|d|+1,|b|+1\}$. Let $\cA$ be the above set of primes. For
large $t$, we have that $2\notin \cA$.

Let $p\in \cA$. Clearly, $2p\ge 2(a-c)t+2(b-d+1)$. Since $a\ge 2c$,
it follows, by an argument similar to the previous one, that the
inequality $2p>at+b$ holds for large values of $t$ unless
$a=2,~c=1$, in which case the inequality $2p>at+b$ also holds for
all primes $p\in \cA$ with $O(1)$ exceptions.

The above argument shows that there exists a positive constant $c_1$
(depending on $a,~b,~c,~d$), such that if $t$ is sufficiently large,
then all primes $p\in \cA$ except for at most $c_1$ of them divide
exactly $x$ (namely, such primes are divisors of the above number
but the squares of these primes are not). Thus, writing $\cB$ for
the subset of $\cA$ such that $p\parallel x$, and putting $u =
\prod_{p \in \cB} p$, it follows that $x = uv$ holds with some
positive integer $v$ coprime to~$u$. Hence,
$$
\sigma(x) = \sigma(u)\sigma(v) = \sigma(v)\prod_{p \in \cB} (p+1).
$$
The factor $p+1$ is even for all $p\in \cB$, and there are
$$
\#\cB \ge \pi(at+b) - \pi((a-c)t+b-d)-c_1
$$
such factors $p+1$. By the Prime Number Theorem, we have that the
inequality $\#\cB\ge ct/(2\log t)$ holds for large values of~$t$.
Hence, writing $\alpha$ for the exponent of the prime $2$ in
$\sigma(x)$, we get that the inequality $\alpha\ge ct/(2\log t)$
holds for large values of~$t$.

By a well-known theorem of Kummer (see, for instance,
Exercise~5.36, p.~245, of the chapter ``Binomial Coefficients''
in~\cite{GKP}),
the exponent $\beta$ of $2$ in
$\binom{n}{m}$ equals the number of carries when adding $m$ and
$n-m$ in base~$2$. In particular, $\beta<c_2 \log t$ for large
values of $t$, where $c_2$ is any fixed constant in $(0,1/\log 2)$.
We are now ready to prove the theorem.

For~(i), let us suppose that $x$ is multiperfect. Then the relation
$$
\sigma(x) = kx
$$
holds for some positive integer~$k$. Letting $\gamma$ be the
exponent of $2$ in the factorization of $k$, we get, by the previous
arguments, that
\begin{equation}
\label{eq:orderof2ink} 
\gamma \ge \alpha-\beta\ge \frac{ct}{2\log t}
- c_2\log t>\frac{ct}{3\log t}
\end{equation}
for large values of~$t$.

On the other hand, since the inequality $x \le 2^n<2^{2at}$ holds
for large values of $t$, we have, by the maximal order of the sum of
divisors function, that
\begin{equation}
\label{eq:cotasup} 
k = \frac{\sigma(x)}{x}\ll \log\log x \le
\log\log(2^{2at}) \le 2\log t
\end{equation}
if $t$ is sufficiently large. Comparing
inequalities~\eqref{eq:orderof2ink} and~\eqref{eq:cotasup}, we get that
$$
2^{ct/(3\log t)} \le 2^{\gamma}\le k\ll \log t,
$$
which shows that $t$ is, in fact, bounded. This completes the proof
of part~(i) of this theorem.

For~(ii), let us suppose that $x$ is $(2,k)$-perfect for some
positive integer~$k$. Then
$$
\sigma(\sigma(x))=kx.
$$
Since $2^{\alpha} \parallel \sigma(x)$, we get that $2^{\alpha+1}-1\mid
\sigma(\sigma(x))$. Recall that if $u\ge 1$ is any positive integer,
then a {\it primitive divisor} of the number $2^u-1$ is a prime
factor of $2^u-1$ which does not divide $2^v-1$ for any positive
integer $1\le v<u$. It is known that a primitive divisor $p$ of
$2^u-1$, whenever it exists, is congruent to $1\pmod u$. It is also
known that if we write
$$
2^u-1=AB,
$$
where $A$ and $B$ are positive integers such that every prime factor
of $A$ is primitive for $2^u-1$ and no prime factor of $B$ is
primitive for $2^u-1$, then  $A>2^{\phi(u)/2}$ holds for all large
values of~$u$. For us, writing
$$
2^{\alpha+1}-1=AB
$$
as described previously, then $A>2^{\phi(\alpha+1)/2}$ provided that
$\alpha$ is large. Since $\alpha>ct/(2\log t)$, and the inequality
$\phi(u)\gg u/\log\log u$ holds for all large positive integers $u$,
it follows that there exists a positive constant $c_3$ such that
\begin{equation}
\label{eq:lowA}
A > \exp\left(\frac{c_3 t}{\log t\log\log t}\right)
\end{equation}
provided that $t$ is sufficiently large. On the other hand, all
prime factors dividing $A$ divide also $kx$. Note that
\begin{equation}
\label{eq:k}
k = \frac{\sigma(\sigma(x))}{x} 
= \frac{\sigma(\sigma(x))}{\sigma(x)} \cdot
\frac{\sigma(x)}{x} \ll (\log\log \sigma(x))\log\log x \ll (\log t)^2,
\end{equation}
where the last inequalities above follow as in part~(i) because
$x<2^{2at}$. If $p\mid A$, then $p\equiv 1\pmod{\alpha+1}$. In
particular, $p>\alpha+1>ct/(2\log t)$, therefore, by inequality
\eqref{eq:k}, we have that $p>k$ for large values of~$t$. Thus,
$p\mid x$. The exponent of $p$ in $x$ is, by Kummer's theorem, at
most $\log t$ for large values of~$t$. Furthermore, $p\le at+b$ and
$p\equiv 1\pmod  {\alpha+1}$, therefore the number of possibilities
for $t$ is at most (even discarding the condition that $p$ is a
prime) $\lfloor (at+b)/(\alpha+1)\rfloor+1<c_4\log t$ for large $t$,
where we can take $c_4=3a/c$. Hence, we have just showed that
\begin{equation}
\label{eq:upA}
A \le \prod_{\substack{p\mid x\\ p\equiv 1\pmod
{\alpha+1}}} p^{\log t}<(at+b)^{c_4(\log t)^2}
< \exp(2c_4(\log t)^3)
\end{equation}
provided that $y$ is sufficiently large. Comparing
inequalities~\eqref{eq:lowA} and~\eqref{eq:upA}, we get that
$$
\frac{c_3 t}{\log t\log\log t} < 2c_4 (\log t)^3,
$$
which obviously implies that $t$ is bounded. This finishes the proof
of part~(ii) of this theorem.



%%%%%%%%%%%%%%%%%%%%%%%%%%%%%%%%%%%%%%%%%
\section{Proof of Theorem~\ref{thm:kge3}}
%%%%%%%%%%%%%%%%%%%%%%%%%%%%%%%%%%%%%%%%%

Assume that $x:=\binom{n}{m}$ is perfect, where $n\ge 2m\ge 6$.
Assume that $n\ge 30$. Corollary~3 in \cite{SaSh} shows that
\begin{equation}
\label{eq:prod}
n(n-1)\cdots (n-m+1)
\end{equation}
is divisible by at least two primes exceeding~$m$. Since $x$ is
perfect, it then follows by a result of Euler that $x=pu^2$ for some
odd prime $p$ and positive integer $u$ (here, $p=2^r-1$ and
$u=2^{(r-1)/2}$ when the perfect number is even). Now $u$ is
divisible by a prime $q$ exceeding~$m$. Hence, $q^2$ divides the
number shown at~\eqref{eq:prod}, and since $q$ exceeds $m$, we get
that $q^2\mid n-i$ for some $i=0,\ldots,m-1$. In particular, $n\ge
(m+1)^2=m^2+2m+1$, therefore $n-m+1>m^2$. Now let $j$ be an integer
in $\{0,1,\ldots,m-1\}$ such that $n-m+1+j$ is a multiple of $p$
(note that $j$ is unique if $p>m$). Since the number shown at
\eqref{eq:prod} is of the form $k!\,pu^2$, it follows easily that if
we remove from this product the number $n-m+1+j$ we get an equation
of the form
$$
(n-m+1)\cdots (n-m+1+(j-1))(n-m+1+(j+1))\cdots n = bw^2,
$$
where $b$ is a positive integer whose largest prime factor is at
most~$m$. This is impossible for $m\ge 3$ by Theorem~2 in~\cite{SaSh}. 
Thus, $n\le 30$, and now a short calculation reveals no
perfect number of the form $\binom{n}{m}$ with $3\le m\le n/2$ and
$n\le 30$.


%%%%%%%%%%%%%%%%%%%%%%%%%%%%%%%%%%%%%%%%%
\section{Proof of Theorem~\ref{thm:2nn}}
%%%%%%%%%%%%%%%%%%%%%%%%%%%%%%%%%%%%%%%%%

We follow the proof of Theorem~\ref{thm:main} but make it
effective. By Corollary~2 (p.~69) in~\cite{RS}, we have
$$
\alpha \ge \pi(2n)-\pi(n) \ge \frac{3n}{5\log n}
$$
for $n\ge 21$. Now, let us observe that $\sigma(\binom{2n}{n})$ is a
multiple of $\prod_{n<p\le 2n} (p+1)$, therefore also of
$2^{\pi(2n)-\pi(n)}$. Letting $\beta$ be the power of $2$ in the
prime factorization of $\binom{2n}{n}$, Kummer's theorem gives us
$2^{\beta}\le 2n$.

We now prove~(i). Assume that $\binom{2n}{n}$ is multiperfect. Let
$k$ be an integer such that $\sigma(\binom{2n}{n}) =
k\binom{2n}{n}$. Then, $2^\gamma\mid k$, where
\begin{equation}
\label{eq:2a} 
2^\gamma \ge 2^{\pi(2n)-\pi(n)-\beta} \ge
\frac{2^{\frac{3n}{5\log n}}}{2n}.
\end{equation}
On the other hand,
$$
k = \frac{\sigma(\binom{2n}{n})}{\binom{2n}{n}} < \sum_{d\le 2^{2n}}
\frac{1}{d} \le 1 + \int_1^{2^{2n}} \frac{dt}{t} \le \log(2^{2n}) +
1 = 2n\log 2 + 1.
$$
Thus, when $n>1$,
$$
\frac{2^{\frac{3n}{5 \log n}}}{2n} \le 2n\log 2 + 1 < 2n,
$$
and then
$$
3 n\log 2 < 10(\log n)\log (2n),
$$
whose largest solution is $n=130$. In this way, we have proved that
$\binom{2n}{n}$ is not multiperfect when $n>130$. For $n \le 130$, a
quick computer search confirmed that $\binom{2n}{n}$ is multiperfect
only when $n=2$.

Assume now that the Catalan number $C_n$ is multiply perfect. Then
the exponent of $2$ in $\sigma(C_n)$ is at least $\pi(2n)-\pi(n)-1$
(by $1$ smaller than before since $n+1$ might be a prime), while the
power of $2$ in $C_n$ does not exceed the power of $2$ in
$\binom{2n}{n}$. Hence, the lower bound on $\gamma$, the power of
$2$ in $k$ defined analogously as $\sigma(C_n)=kC_n$, satisfies the
inequality
\begin{equation}
\label{eq:2aCatalan} 
2^\gamma \ge \frac{2^{\frac{3n}{5\log n}}}{4n}
> \frac{2^{\frac{3n}{5\log n}}}{(2n)^2} .
\end{equation}
Thus, following the same arguments as before, we arrive at
$$
3 n\log 2 < 15(\log n)\log (2n),
$$
whose largest solution is $n=247$. A quick computer check confirmed
that the only $n\le 247$ for which $C_n$ is multiply perfect is
$n=1$.

We now look at~(ii). Let $x=B_n$ or $C_n$ and let $\alpha$ be the
power of $2$ in either $\sigma(B_n)$ or $\sigma(C_n)$. Then the
arguments from part~(i) show that
$$
\alpha\ge \pi(2n)-\pi(n)-1\ge \frac{3n}{5\log n}-1
\qquad {\text{for}}~n\ge 1.
$$
Now $2^{\alpha+1}-1$ divides $\sigma(\sigma(x))=kx$. Clearly,
$\sigma(x)<x^2<2^{4n}$, so the arguments from~(i) show that
\begin{align*}
k & = \frac{\sigma(\sigma(x))}{\sigma(x)} \cdot \frac{\sigma(x)}{x}
< \left(1+\int_1^{2^{2n}} \frac{dt}{t}\right)
\left(1+\int_1^{2^{4n}} \frac{dt}{t}\right) \\
& <  (1+2n\log 2)(1+2n\log 4) < (2n)^2,
\end{align*}
where the last inequality holds for $n>30$, which we are assuming
from now on. Now write
$$
2^{\alpha+1}-1=AB,
$$
where all prime factors of $A$ are congruent to $1$ modulo
$\alpha+1$. It is well-known, from the proof of the Primitive
Divisor Theorem (see, for example, \cite{Z}), that
$$
A > \frac{2^{\phi(\alpha+1) - 2^{\omega(\alpha+1)-1}}}{\alpha+1},
$$
where for a positive integer $m$ we put $\omega(m)$ for the number
of distinct prime factors of~$m$. We now assume that $n>135$. Then
$$
\alpha+1 > \frac{3n}{5\log n} > 12
$$
(where the last inequality holds for $n>90$), therefore, by a result
of Schinzel, \cite{Sch}, we have that $A$ has a prime factor at
least as large as $2(\alpha+1)+1$. This prime factor either divides
$x$ (so it is $<2n$), or it divides $k<4n^2$, therefore
$\alpha+1<2n^2$. Thus,
\begin{equation}
\label{eq:Alow}
A > \frac{2^{\phi(\alpha+1)}-2^{\omega(\alpha+1)-1}}{2n^2}.
\end{equation}
Now every prime factor dividing $A$, either divides $k$, or divides~$x$.
Let $p$ be a prime factor dividing both $A$ and~$x$. Then
$$
p^2 > (\alpha+1)^2 > \left(\frac{3 n}{5\log n}\right)^2 > 2n,
$$
where the last inequality holds for all $n>135$, which we are
assuming. Thus, $p \parallel x$. Furthermore, the number of such primes is
$\le 2n/(\alpha+1)$ and every one of such primes is~$<2n$. This
shows that
$$
A\le k(2n)^{2n/(\alpha+1)}<(2n)^{2/(\alpha+1)+2}.
$$
Thus, comparing the above upper bound on $A$ with the lower bound
\eqref{eq:Alow}, we get
\begin{equation}
\label{eq:master}
\phi(\alpha+1) - 2^{\omega(\alpha+1)-1} 
< \left(\frac{2n}{\alpha+1}+4\right)
\frac{\log(2n)}{\log 2}.
\end{equation}
This is the master inequality. Since $\alpha+1>3n/(5\log n)$, we get
that
\begin{equation}
\label{eq:up}
\phi(\alpha+1)-2^{\omega(\alpha+1)-1} 
< \left(\frac{10\log n}{3}+4\right) \frac{\log(2n)}{\log 2}.
\end{equation}
Let $s$ be the number of prime factors of $\alpha+1$. Clearly,
$2^s\le \alpha+1$, so $s\le (\log(\alpha+1))/\log 2$, therefore
$$
\phi(\alpha+1) \ge
(\alpha+1) \prod_{i=2}^{s+1}\left(1-\frac{1}{i}\right) \ge
\frac{\alpha+1}{s+1}\ge \frac{(\alpha+1)\log 2}{\log(2(\alpha+1))}.
$$
Finally, writing $\tau(m)$ for the number of divisors of $m$, and
observing that
$$
2^{\omega(\alpha+1)-1}\le \frac{\tau(\alpha+1)}{2} \le
\frac{\sqrt{3(\alpha+1)}}{2},
$$
we get that
$$
\phi(\alpha+1)-2^{\omega(\alpha+1)-1}
\ge \frac{(\alpha+1)\log 2}{\log(2(\alpha+1))}
- \frac{\sqrt{3(\alpha+1)}}{2}.
$$
The function
$$
t\mapsto \frac{t\log 2}{\log(2t)} - \frac{\sqrt{3t}}{2}
$$
is increasing for $t\ge 8$, and since for us 
$\alpha+1\ge 3n/(5\log n)>12$, we have that
\begin{equation}
\label{eq:down} 
\phi(\alpha+1)-2^{\omega(\alpha+1)-1} \ge
\frac{3n\log 2}{5(\log n)\log(6n/(5\log n))}
- \frac{3}{2}\left(\frac{n\log 2}{5}\right)^{1/2}.
\end{equation}
From inequalities~\eqref{eq:up} and~\eqref{eq:down}, we get that
$$
\frac{3n\log 2}{5(\log n)\log(6n/(5\log n))}
- \frac{3}{2}\left(\frac{n\log 2}{5\log n}\right)^{1/2}
< \left(\frac{10\log n}{3}+4\right) \frac{\log(2n)}{\log 2},
$$
giving $n\le 310000$. Thus, there is no solution with $n>310000$.

Now we start lowering the upper bound on~$n$. Assume  still $n>135$
but $n\le 310000$. Here, we have that $\alpha+1\le \pi(2n)+1<51000$.
The product of the first $7$ primes exceeds $51000$, therefore
$$
\frac{\phi(\alpha+1)}{\alpha+1}\ge \prod_{2\le q\le 13}
\left(1-\frac{1}{q}\right)=\frac{192}{1001}.
$$
Moreover, for $\alpha+1<51000$, we have $\tau(\alpha+1)\le 108$,
therefore the largest power of $2$ which is $\le \tau(\alpha+1)/2$
is at most~$32$. Hence, we get the inequality
$$
\phi(\alpha+1)-2^{\omega(\alpha+1)-1} \ge
\frac{192(\alpha+1)}{1001}-32\ge \frac{594n}{5005\log n}.
$$
Comparing the last inequality above with inequality~\eqref{eq:up},
we get
$$
\frac{594 n}{5005 \log n}-32 
< \left(\frac{10\log n}{\log 3}+4\right)
\frac{\log(2n)}{\log 2},
$$
giving $n\le 72000$. Thus, there is no solution with $n>72000$.
Assume now that $n>20000$. We have checked with Mathematica that
$\pi(2n)-\pi(n)>0.95 n/\log n$ for all $n\in [20000,80000]$. Thus,
instead of the inequality $\alpha+1>3 n/(5\log n)$, we can use the
inequality $\alpha+1>0.95 n/\log n$, in which case we also have
$2n/(\alpha+1)<40(\log n)/19$. Furthermore, since $2n<144000$ and
$\alpha+1<\pi(144000)<14000$, but the product of the first $6$
primes exceeds $14000$, it follows that
$$
\frac{\phi(\alpha+1)}{\alpha+1}\ge \prod_{2\le q\le 11}
\left(1-\frac{1}{q}\right)=\frac{16}{77}.
$$
Thus,
$$
\phi(\alpha+1)-2^{\omega(\alpha+1)-1}
\ge \frac{15.2 n}{77\log n}-32,
$$
and comparing this with the bound~\eqref{eq:master}, we get
$$
\frac{15.2 n}{5\log n}-32 < \left(\frac{40\log n}{19}
+ 4\right) \frac{\log(2n)}{\log 2},
$$
giving $n<20000$. Thus, there is no solution with $n\ge 20000$.

Next we checked that there is no solution with $n>1100$ in the
following way. For each $n\in [1100,20000]$, we used Mathematica to
find the exact value of~$\alpha$. In order to do so, for each $n$ we
factored the integer $B_n$ and $C_n$ (which are large integers all
whose prime factors are smaller than $2n$), and then
$$
\alpha = \sum_{\substack{p^{\alpha_p}\mid B_n\\
\alpha_p~{\text{odd}}}} \operatorname{ord}_2
\left(\frac{p^{\alpha_p+1}-1}{p-1}\right),
$$
where for a positive integer $m$ we write $\operatorname{ord}_2(m)$
for the exponent of $2$ in the factorization of~$m$. We then checked
whether the master inequality~\eqref{eq:master} holds. No solution
for either $x=B_n$ or $x=C_n$ was found with $n>1100$. We then
checked that the only $n\le 1100$ for which $x \mid
\sigma(\sigma(x))$ are the claimed ones.


%%%%%%%%%%%%%%%%%%%%%%%%%%%%%%%%%%%%%%%%%
\section{Proof of Theorem~\ref{thm:dens}}
%%%%%%%%%%%%%%%%%%%%%%%%%%%%%%%%%%%%%%%%%

We start with part~(i) and assume, as in the proof of
Theorem~\ref{thm:main}, that we are looking at pairs $(n,m)$
parameterized as in Lemma~\ref{lem:param}. Let $t$ be large and put
$z_1=\exp((\log t)^{1/2})$ and $z_2=t^{1/4}$. Lemma 5.1 with $J=1$
and $\varepsilon=1/10$ in~\cite{GrR} shows that for any two
constants $0<\sigma_0<\sigma_1<1$ there exists some positive
constant $c_1$ such that uniformly in $y\le x$ we have that the set
of prime numbers $p\le y$ such that $\{x/p\} \in (\sigma_0,\sigma_1)$
has cardinality
$$
(\sigma_1-\sigma_0)\pi(y) + O\left((y^{1-c_1(\log y)^2/\log x}
+ y^{8/5}x^{-1/2})(\log x)^4\right).
$$
From the way we have chosen $z_1$ and $z_2$, we deduct that for
large $t$, there are
$$
(\sigma_1-\sigma_0) \pi(y) \left(1+O\left(\frac{1}{\log y}\right)\right)
$$
primes $p\le y$ such that $\{t/p\}\in (\sigma_0,\sigma_1)$ and this
holds for all $y\in [z_1,z_2]$. We now put
$\sigma_0=1-1/(4c(a-c)),~\sigma_1=1-1/(5c(a-c))$, and deduct that if
we put $\cP_y$ for the set of primes $p\in (y,2y)$ with $\{t/p\}$ in
the above interval, then the inequality
$$
\# \cP_y\ge \frac{y}{30c(a-c)\log y}
$$
holds uniformly for all $y\in [z_1,z_2/2]$ once $t$ is sufficiently
large. Since all such primes are at most as large as $y$, we get
that
\begin{equation}
\label{eq:sum1overp} 
\sum_{p\in \cP_y}\frac{1}{p} \ge
\frac{c_1}{\log y},
\end{equation}
where we put $c_1=1/(30c(a-c))$. Observe that if $p$ is such a
prime, then
$$
\left\{\frac{(a-c)t+(b-d)}{p}\right\} \in \left(1-\frac{1}{3c},
1-\frac{1}{6c}\right),
$$
and
$$
\left\{\frac{ct+d}{p}\right\} \in \left(1-\frac{1}{3(a-c)},
1-\frac{1}{6(a-c)}\right),
$$
provided that $t$ is large. In particular, since
$$
\left(1-\frac{1}{3c}\right) + \left(1-\frac{1}{3(a-c)}\right) > 1
$$
for all $c\ge 1$ and $a-c\ge 1$, we get that for large $t$, the
binomial coefficient $x=\binom{n}{m}$ is a multiple of such a prime
$p$ since there is at least one carry when adding $(a-c)t+(b-d)$
with $ct+d$ in base~$p$.

Now give $y$ the values $2^{\ell}$, where $\ell$ is an integer.
Since $2^{\ell}$ must belong to the interval $[z_1,z_2/2]$, it
follows that we can take $\ell\in [L_1,L_2]$, where
$$
L_1 = \left\lfloor \frac{(\log t)^{1/2}}{\log 2}\right\rfloor + 1
\qquad {\text{and}}\qquad L_2
= \left\lfloor \frac{\log t}{4\log 2}\right\rfloor-1.
$$
Thus, $x$ is divisible by all primes in $\bigcup_{L_1\le \ell
\le L_2} \cP_{2^{\ell}}$, and the sets making up the above
union are disjoint. Hence, using~\eqref{eq:sum1overp}, we get
\begin{align*}
\frac{\sigma\left(x\right)}{x}
& \ge \prod_{L_1\le \ell\le L_2}
\prod_{p\in \cP_{2^{\ell}}}\left(1+\frac{1}{p}\right)
= \exp\left(\sum_{\substack{L_1\le \ell\le L_2\\ p\in \cP_{2^{\ell}}}}
\frac{1}{p} + O\left(\sum_{p\ge 2} \frac{1}{p^2}\right)\right)\\
& \ge \exp\left(c_2\sum_{L_1\le \ell\le L_2} \frac{1}{\ell} +
O(1)\right)
\gg \exp\left(c_2\int_{L_1}^{L_2}\frac{ds}{s}\right) 
\gg \exp\left(c_2\log\frac{L_1}{L_2}\right) \gg (\log t)^{c_L},
\end{align*}
where we put $c_2=c_1/\log 2$ and $c_L=c_2/2$. Since $x\le 2^{2at}$
for large values of $t$, we get that
$$
\frac{\sigma(x)}{x} \gg \left(\log\log x\right)^{c_L},
$$
where $c_L$ depends only on the line $L$, which is what we wanted to
prove.

Finally, for~(ii) let $m$ be fixed, $\alpha\in [1,\infty)$, and
$\varepsilon>0$ be arbitrary. Choose $m$ disjoint finite sets of
primes $\cP_i$ for $i=1,\ldots,m$ all exceeding $m$ such
that
$$
\prod_{p\in \cP_i}\left(1+\frac{1}{p}\right)\in
[\alpha^{1/m}(1-\varepsilon),\alpha^{1/m}(1+\varepsilon)].
$$
Let also $t$ be some positive real number whose natural logarithm
exceeds all members of $\bigcup_{1\le i\le m} \cP_i$ and
let $\cQ$ be the set of primes $p\in (m,\log t)$ not
belonging to any of the sets $\cP_i$ for $i=1,\ldots,m$.
Put $M=(m!)^2$, $P_i=\prod_{p\in \cP_i} p$,
$P=\prod_{i=1}^m P_i$, $Q=\prod_{q\in \cQ} q$ and consider
the following system of congruences
$$
n\equiv 0\pmod {MQ}\qquad \text{and}\qquad n+i\equiv P_i\pmod
{P_i^2} \quad \text{for}~i=1,\ldots,m.
$$
By the Chinese Remainder Theorem, the above system of congruences
puts $n$ into an arithmetic progression $n\equiv n_0\pmod {N}$,
where $N=MQP^2$. The size of the modulus is, by the Prime Number
Theorem,
$$
N\le M\left(\prod_{p\le \log t} p\right)^2=Me^{2(1+o(1))\log t}
$$
as $t\to\infty$, therefore $N<t^3$ for large values of~$t$. Hence,
there is a number $n\in [t^3,2t^3)$ in the above progression once
$t$ is sufficiently large. Observe that
$$
x=\binom{n+m}{m}=\prod_{i=1}^m\left(1+\frac{n}{i}\right).
$$
Putting $m_i=1+n/i$, it is easy to see that $m_i\equiv 1\pmod {m!}$,
that $m_i$ is divisible by all primes $p$ in $\cP_i$ but
not by their squares, and $m_i$ is not divisible by any member of
$\cQ$ for any $i=1,\ldots,m$. Hence, $m_i=P_iM_i$, where
$M_i$ is a positive integer whose prime factors exceed $\log t$.
Since $M_i<2t^3+1$, we get that $\omega(M_i)=O(\log t/\log\log t)$
for all $i=1,\ldots,m$. Thus,
$$
x=\binom{n+m}{m} = \prod_{i=1}^m P_i\prod_{i=1}^m M_i,
$$
where the two products above are coprime. Thus,
$$
\frac{\sigma(x)}{x} 
= \frac{\sigma\left(\prod_{i=1}^m P_i\right)}{\prod_{i=1}^m P_i}
\, \frac{\sigma\left(\prod_{i=1}^m M_i\right)}{\prod_{i=1}^m M_i}.
$$
Note that, from the way we have chosen our sets $\cP_i$, we
have
$$
\frac{\sigma\left(\prod_{i=1}^m P_i\right)}{\prod_{i=1}^m P_i}
= \prod_{i=1}^m \prod_{p\in \cP_i}
\left(1+\frac{1}{p}\right) \in
[\alpha(1-\varepsilon)^m,\alpha(1+\varepsilon)^m],
$$
while
$$
1\le \frac{\sigma\left(\prod_{i=1}^m M_i\right)}{\prod_{i=1}^m M_i}
\le \left(1+\frac{1}{\log t}\right)^{O(m\log t/\log\log t)}
= 1+O\left(\frac{m}{\log\log t}\right),
$$
so this last number is in $[1,1+\varepsilon)$ for sufficiently large
values of~$t$. Since $\alpha$ and $\varepsilon$ were arbitrary and
$m$ was fixed, we get that $\alpha$ can be approximated arbitrarily
well by rational numbers of the form $\sigma(x)/x$ with
$x=\binom{n+m}{m}\in L$, which is what we wanted to prove.


\paragraph*{Acknowledgements.} We thank Professor Igor Shparlinski for a
useful discussion. We also thank the anonymous referee and Professor
Carl Pomerance for useful suggestions.


%   --- --- --- --- --- ---
%
%   \url{http://www.research.att.com/~njas/sequences/A007691}
%
%   \url{http://en.wikipedia.org/wiki/Multiply_perfect_number}
%
%   \url{http://mathworld.wolfram.com/MultiperfectNumber.html}
%
%   \url{http://wwwhomes.uni-bielefeld.de/achim/mpn.html}
%
%   --- --- --- --- --- ---


%%%%%%%%%%%%%%%%%%%%%%%%%%%%%%%%%%%%%%%%%
\begin{thebibliography}{99}
%%%%%%%%%%%%%%%%%%%%%%%%%%%%%%%%%%%%%%%%%

\bibitem{D} L. E. Dickson, 
\textit{History of the Theory of Numbers}, vol.~1,
Carnegie Institute of Whashington, 1919. Reprinted by Chelsea,
1971.

\bibitem{Ev} U. Everling, Solution to ``Problem 10331'',
\textit{Amer. Math. Monthly} \textbf{103} (1996), 701--702.

\bibitem{GKP} R.~L.~Graham, D.~E.~Knuth and O.~Patashnik,
% "Binomial Coefficients." Ch. 5 (pp. 153--242) in
\textit{Concrete Mathematics: A Foundation for Computer Science},
Addison-Wesley, Reading, MA, 1989.

\bibitem{GrR} A.~Granville and O.~Ramar\'e, Explicit bounds
on exponential sums and the scarcity of squarefree binomial
coefficients, \textit{Mathematika} \textbf{43} (1996), 73--107.

\bibitem{Guy} R.~K.~Guy, \textit{Unsolved Problems in Number Theory},
third ed., Springer, 2004.

\bibitem{Leh} D.~N.~Lehmer, Multiply perfect numbers,
\textit{Ann. of Math. (2)} \textbf{2}
(1900/1901), 103--104.

\bibitem{Lu1} F.~Luca, Perfect Fibonacci and Lucas numbers, \textit{Rend. Circ.
Mat. Palermo (2)} \textbf{49} (2000), 313--318.

\bibitem{Lu2} F.~Luca, The anti-social Fermat number, \textit{Amer. Math.
Monthly} \textbf{107} (2000), 171--173.

\bibitem{Lu3} F.~Luca, Multiply perfect numbers in Lucas sequences with
odd parameters, \textit{Publ. Math. Debrecen} \textbf{58} (2001),
121--155.

\bibitem{Lu4P} F.~Luca, Problem H-599, % (Catalan Fibonacci numbers)
\textit{Fibonacci Quart.} \textbf{41} (2003), 380.

%\bibitem{Lu4S} F.~Luca, Solution to ``Problem H-599'',
%\textit{Fibonacci Quart.} \textbf{?}, ???

\bibitem{Lu5} F.~Luca, On the sum of divisors of the Mersenne numbers,
\textit{Math. Slovaca} \textbf{53} (2003),  457--466.

\bibitem{Lu6} F.~Luca, Sums of divisors of
binomial coefficients, \textit{Int. Math. Res. Not. IMRN}
\textbf{2007}, Article ID rnm101.

\bibitem{L-P} F.~Luca, Problem 73,
\textit{Gac. R. Soc. Mat. Esp.} \textbf{10} (2007), 145 (in
Spanish).

\bibitem{L-S} F.~Luca, Solution to ``Problem 73'',
\textit{Gac. R. Soc. Mat. Esp.} \textbf{11} (2008), 98--99 (in
Spanish).

\bibitem{LuPa} F.~Luca and F.~Pappalardi, Members of binary recurrences
on lines of the Pascal triangle, \textit{Publ. Math. Debrecen}
\textbf{67} (2005), 103--113.

\bibitem{LuPo} F.~Luca and C.~Pomerance,
On some problems of M\c{a}kowski-Schinzel and Erd\H{o}s concerning the
arithmetical functions $\phi$ and~$\sigma$,
\textit{Colloq. Math.} \textbf{92} (2002),
111--130.

\bibitem{LuSh} F.~Luca and I.~E.~Shparlinski, Arithmetic functions with
linear recurrence sequences, \textit{J. Number Theory} \textbf{125}
(2007), 459--472.

\bibitem{Po} C. Pomerance, Problem 10331,
\textit{Amer. Math. Monthly} \textbf{100} (1993), 796.

\bibitem{RS} J.~B.~Rosser and L.~Schoenfeld, Approximate formulas
for some functions of prime numbers, \textit{Illinois J. Math.}
\textbf{6} (1962), 64--94.

\bibitem{SaSh} N.~Saradha and T.~N.~Shorey, Almost squares and factorisations in
consecutive integers, \textit{Compositio Math.}  \textbf{138}
(2003), 113--124.

\bibitem{Sch} A.~Schinzel, On primitive prime factors of $a^n-b^n$,
\textit{Proc. Cambridge Philos. Soc.} \textbf{58} (1962), 555--562.

\bibitem{Shp} I.~E.~Shparlinski,
Some arithmetic properties of recurrence sequences, \textit{Mat.
Zametki} (\textit{Math. Notes}) \textbf{47} (1990), 612--617 (in Russian).

\bibitem{njas}
N. J. A. Sloane, Multiply-perfect numbers, \textit{The On-Line
Encyclopedia of Integer Sequences},
\url{http://www.research.att.com/~njas/sequences/A007691}

\bibitem{Z} K.~Zsigmondy, Z\"ur Theorie der Potenzreste, \textit{Monatsh. Math.}
\textbf{3} (1892), 265--284.

\end{thebibliography}

\end{document}
